\renewcommand{\bibname}{Sources and authority map}
\begin{thebibliography}{99}
\addcontentsline{toc}{chapter}{Sources and authority map}
\raggedright
\bibitem{onsager-one} Lars Onsager. \emph{Reciprocal Relations in Irreversible Processes. I.} Physical Review 37, 405--426 (1931). \href{https://doi.org/10.1103/PhysRev.37.405}{doi:10.1103/PhysRev.37.405}. Historical primary source; no economic application is inferred.
\bibitem{onsager-two} Lars Onsager. \emph{Reciprocal Relations in Irreversible Processes. II.} Physical Review 38, 2265--2279 (1931). \href{https://doi.org/10.1103/PhysRev.38.2265}{doi:10.1103/PhysRev.38.2265}. Linear response, microscopic reversibility and dissipation-function context. The EBU threshold-mobility controller is a separately declared adaptation.
\bibitem{programme} Energy Balance Project. \emph{Local Gaussian EBU: Programme Reconciliation}. Committed source \texttt{LOCAL\_GAUSSIAN\_EBU\_PROGRAMME\_RECONCILIATION.md}, inspected at \texttt{05f3cb0}. Sections 3--8 supply the prospective domain and conditional accounting identities; Section 10 retains the open decisions. Not an executed study.
\bibitem{foundation} Energy Balance Project. \emph{V3.0 Local EBU Foundation Draft}. \texttt{V3.0\_LOCAL\_EBU\_FOUNDATION\_DRAFT.md}. Sections 6--8: exact finite quotation, cost categories, local cancellation and scoped sequential identities. Historical D0/P1C assumptions remain attached to their statements.
\bibitem{bridge} Energy Balance Project. \emph{The Sequential--Parallel Bridge of EBU}, analytical checkpoint v0.2. \texttt{SEQUENTIAL\_PARALLEL\_BRIDGE.md}, Sections 3--7. Joint state changes, telescoping and comparator-relative interaction. Its operational status passages are historical, not a current Gate 1D-C completion record.
\bibitem{framework} Energy Balance Project. \emph{Local Gaussian Programme: Framework Decision}. \texttt{LOCAL\_GAUSSIAN\_EBU\_FRAMEWORK\_DECISION.md}. Source-locked F2 reuse decision for planning, with separate event, ownership, capacity and locality dependencies. Does not authorize scientific runtime.
\bibitem{roadmap} Energy Balance Project. \emph{Local Gaussian EBU: Deferred Implementation Roadmap}. \texttt{LOCAL\_GAUSSIAN\_EBU\_IMPLEMENTATION\_ROADMAP.md}. Implementation, transition fixtures, preregistration and execution are separate stages. Historical book-stage wording is superseded only by the author's later book-revision instruction.
\bibitem{gate} Energy Balance Project. \emph{Gate 1D-C Finalized Execution Manifest}, \texttt{results/v3.0/gate1dc/MANIFEST.md}; accompanying registered protocol, JSON plan, summary and stdout. Historical finalized evidence: 30 runs and 6,000 trace rows. Execution identity \texttt{8793bda}. No Gaussian inference is made.
\bibitem{originals} Konrad Grzyb. \emph{The Energy Balance Project: Unified Explanatory Edition}, Parts II and III. Author-supplied PDFs, 160 and 153 pages. Content/style baselines and sources of the explicitly marked historical appendices. Original artifacts remain unchanged.
\bibitem{partone} Konrad Grzyb. \emph{Physical Intuition and the EBU Idea}, revised Gaussian explanatory Part I. Author-review edition. The shared Ridge--Vale examples and typographic language continue here; examples are mathematical teaching constructions, not experimental evidence.
\end{thebibliography}
\section*{How to verify a source claim}
The companion source manifest records full file hashes, the checkout identity,
the original PDF identities and the exact original pages retained in each
appendix. The outline disposition records subjects rewritten rather than
claiming every original sentence survives. All new derivations state their
assumptions in the chapter itself. A citation identifies a source; it does
not transform a prospective design into an accepted theorem or experiment.
